% Bewertungspaket zum Gesamttest «Vierecke» — Quelle des PDFs (python3 scripts/build-lp-pdf.py vierecke).
% Ganze Punkte je Teilschritt; Werte mit python3 nachgerechnet (scripts/lp/vierecke/zahlen.py, 08.10.2026,
% Folgefehler-Fälle eingeschlossen). Aufbau wie trigonometrische-berechnungen/bewertungspaket.tex.
\documentclass[11pt]{article}
\usepackage{../lp-druck}
\renewcommand{\lpname}{Vierecke}
\renewcommand{\lpbereich}{Grundlagenfach}
\renewcommand{\lpteilgebiet}{5.2b}
\renewcommand{\lpfuss}{Bewertungspaket}
\newcolumntype{P}{>{\centering\arraybackslash}p{10mm}}
\newenvironment{raster}{\par\smallskip\noindent\tabularx{\linewidth}{@{}>{\raggedright\arraybackslash}p{38mm}P X@{}}\toprule
  \small\textsc{Teilschritt} & \small\textsc{P} & \small\textsc{Musterlösung}\\\midrule}{\bottomrule\endtabularx\par}
\newcommand{\fehler}[1]{\par\smallskip{\small\textcolor{lprot}{\bfseries Typische Fehler:} #1}}
\newcommand{\E}{\,{\footnotesize\bfseries\color{lpgruen}(E)}}
\newcommand{\F}{\,{\footnotesize\bfseries\color{lporange}(F)}}
\newcommand{\A}{\,{\footnotesize\bfseries\color{lpblau}(A)}}
\begin{document}
\lpkopf{Bewertungspaket zum Gesamttest}{}

\begin{lpachtung}{Erst lösen, dann öffnen}
Dieses Paket enthält die Musterlösung. Wer es vor dem Lösen liest, misst am Ende nur, wie gut das Abschreiben gelingt.
\end{lpachtung}

\section*{1 · So gehst du vor}
\begin{enumerate}[leftmargin=*,itemsep=0pt]
\item \textbf{Gesamttest lösen} — vollständig, mit Skizze und Rechenweg, ohne dieses Paket; der Taschenrechner ist erlaubt.
\item \textbf{Fotografieren oder scannen}, gut lesbar, eine Seite pro Bild. \textbf{Kein Name} auf den Bildern.
  Handyfotos tragen oft unsichtbar Standort, Zeit und Gerät mit (Metadaten). Lade darum lieber einen Scan oder ein
  Bildschirmfoto des Fotos hoch, oder entferne vorher die Standortangaben.
\item \textbf{Bewerten lassen.} Ob du dafür eine KI nutzt und welche, entscheidest du selbst und in eigener Verantwortung —
  je nachdem, was dir aufgrund deines Alters und deiner persönlichen Situation erlaubt ist (Altersgrenzen und
  Nutzungsbedingungen des Anbieters, allenfalls Einverständnis deiner Eltern). Lade nur deine Lösung und dieses PDF hoch,
  keine persönlichen Daten, und füge den Auftrag aus Abschnitt~2 ein. Ohne KI kannst du dich mit Abschnitt~3 selbst bewerten.
\item \textbf{Rückmeldung prüfen.} Stimmt, was die KI aus deiner Handschrift gelesen hat? Eine KI kann sich irren —
  im Zweifel gilt das Raster in Abschnitt~3.
\item \textbf{Gezielt wiederholen} nach Abschnitt~4.
\end{enumerate}

\needspace{8\baselineskip}
\section*{2 · Auftrag an die KI}
\begin{tcolorbox}[breakable,colback=lplinie!25,colframe=lplinie,boxrule=0.5pt,arc=1mm,fontupper=\small,left=2mm,right=2mm,top=1.5mm,bottom=1.5mm]
Du bist Korrektorin oder Korrektor für Mathematik an einer Schweizer Berufsmaturitätsschule. Im Anhang findest du (1)~das Bewertungspaket zum Gesamttest «Vierecke» mit Musterlösung und Punkteraster und (2)~Fotos meiner handschriftlichen Lösung.\par\smallskip
Geh so vor:\par
1. Schreib zu jeder Aufgabe G1 bis G6 zuerst ab, was du in meiner Lösung liest — Rechenweg und Ergebnis. Was du nicht sicher lesen kannst, markierst du mit [unsicher]. Nicht raten.\par
2. Vergib die Punkte genau nach dem Raster im Bewertungspaket (Abschnitt 3), ganze Punkte je Zeile. Die Spalte P sagt, wie eine Zeile zählt: \textbf{(E)} — das Ergebnis muss stimmen und aus einem erkennbaren Weg hervorgehen. \textbf{(F)} — Folgezeile: Sie zählt auch, wenn ich mit einem eigenen, falschen Vorwert richtig weitergerechnet habe, ausser dieser Vorwert ist \emph{unmöglich} (Regel in Abschnitt 3). \textbf{(A)} — erkennen oder benennen, zählt auch ohne Rechenweg. Ohne Zeichen — die Zeile bewertet den Weg; auch sie zählt mit einem eigenen, möglichen Vorwert. So wird jeder Fehler nur einmal abgezogen. Die Zeilen «Typische Fehler» sagen, welche Rasterzeilen bei diesem Fehler 0 Punkte geben — sie sind kein zusätzlicher Abzug.\par
3. Andere richtige Lösungswege zählen voll, ebenso gleichwertige Schreibweisen: Dezimalpunkt oder Dezimalkomma, exakte Werte wie $\sqrt{32.5}$ statt Dezimalzahlen, anders, aber vernünftig gerundete Ergebnisse ($\pm 0.01$ oder durch gerundete Zwischenwerte erklärbar). \textbf{Einheiten:} Fehlt bei einem Endergebnis die Einheit oder ist sie falsch ($\text{cm}$ statt $\text{cm}^2$), gibt diese Zeile 0 Punkte — aber nur beim ersten Mal im ganzen Gesamttest; weitere Einheitenfehler kosten nichts mehr.\par
4. Begründe jeden Abzug in einem Satz und nenne die Fehlerart (zum Beispiel «plus statt minus im Pythagoras» oder «ganzer Unterschied statt Überstand»). Schreib die Musterlösung nicht ab — zeig mir, wo mein Fehler liegt.\par
5. Gib zum Schluss eine Tabelle «Aufgabe | Punkte | Begründung», die Gesamtpunktzahl (von 24) und — nach der Selbsteinschätzung in Abschnitt 4 — welches Kapitel des Leitprogramms ich wiederholen soll. Keine Note.\par
6. Fehlt ein Foto oder ist eine Aufgabe unlesbar, sag es, statt sie zu bewerten.
\end{tcolorbox}

\needspace{8\baselineskip}
\section*{3 · Musterlösung und Punkteraster}
\textbf{Allgemein:} Ganze Punkte je Zeile. In der Spalte P: \textbf{(E)}~= Ergebnis, das nur aus den gegebenen Zahlen
folgt — richtig \emph{und} Weg erkennbar · \textbf{(F)}~= Folgezeile: zählt auch mit einem eigenen, falschen Vorwert,
wenn damit richtig gerechnet ist · \textbf{(A)}~= erkennen oder benennen, zählt auch ohne Rechenweg · ohne Zeichen = die
Zeile bewertet den Weg (auch mit einem eigenen Vorwert). So wird jeder Fehler nur einmal abgezogen.
\textbf{Unmöglich} ist ein eigener Vorwert, der der Figur widerspricht: eine Länge, die negativ oder null ist; eine
Kathete (Höhe, Seite, halbe Diagonale), die länger ist als die Hypotenuse im selben Teildreieck; eine Mittellinie, die
nicht zwischen $a$ und $c$ liegt; ein Viereckswinkel ab $180^\circ$. Auf einem unmöglichen Vorwert gibt es keine Folgepunkte.
Taschenrechner erlaubt; Längen und Flächen auf zwei Dezimalen, Toleranz $\pm 0.01$ (oder durch gerundete Zwischenwerte
erklärbar); Folgewerte aus den ungerundeten eigenen Werten. \textbf{Einheiten:} Der erste Einheitenfehler im ganzen
Test kostet die Zeile dieses Ergebnisses, weitere nichts. «Typische Fehler» nennt, welche Zeilen 0 Punkte geben, und
die Punkte, die bleiben. \textbf{Teilrichtige Zeilen:} Verlangt eine Zeile mehrere Angaben und stimmt nur ein Teil,
gibt die Zeile 0 Punkte.

\begin{lpaufgabe}{G1}{Vierecke beschreiben}{4 P}
\begin{raster}
(a) Viereck & 1\A & Ein \textbf{Rhombus} (Raute). «Rhombus, kein Quadrat» ist gleichwertig.\\
(a) Begründung & 1 & Die Diagonalen halbieren sich: Parallelogramm. Dazu stehen sie senkrecht: Rhombus. Sie sind verschieden lang ($7 \neq 10$): kein Rechteck, also kein Quadrat. Nötig sind alle drei Schritte.\\
(b) $\alpha$, $\delta$ & 1\E & $\alpha$ und $\delta$ liegen am Schenkel $d$: $\alpha + \delta = 180^\circ$ und $\delta - \alpha = 40^\circ$, also $\alpha = 70^\circ$, $\delta = 110^\circ$. Beide nötig.\\
(b) $\gamma$ & 1\E & $\beta$ und $\gamma$ liegen am Schenkel $b$: $\gamma = 180^\circ - 80^\circ = 100^\circ$. (Probe: $70^\circ + 80^\circ + 100^\circ + 110^\circ = 360^\circ$.)\\
\end{raster}
\fehler{«Quadrat»: beide (a)-Zeilen 0 (die Diagonalen sind verschieden lang) $\to$ höchstens 2 von 4\,P. «Parallelogramm» oder «Drachen»: (a)-Viereck 0; die Begründung zählt nur mit allen drei Schritten.
Begründung ohne «halbieren» (nur «senkrecht $\Rightarrow$ Rhombus»): Begründungszeile 0.
$\alpha + \delta = 360^\circ$ gerechnet ($\alpha = 160^\circ$, $\delta = 200^\circ$): $\alpha$-$\delta$-Zeile 0.
$\alpha$ mit $\beta$ statt mit $\delta$ gepaart ($\alpha = 100^\circ$, $\delta = 140^\circ$): $\alpha$-$\delta$-Zeile 0.
$\gamma = \alpha$ oder $\gamma = 80^\circ$ (gegenüberliegende Winkel gleich gesetzt wie im Parallelogramm): $\gamma$-Zeile 0.}
\end{lpaufgabe}

\begin{lpaufgabe}{G2}{Parallelogramm}{4 P}
\begin{raster}
(a) Teildreieck & 1 & Das Dreieck $AFD$ ist rechtwinklig bei $F$, die Seite $\overline{AD}$ ist seine Hypotenuse: $h = \sqrt{\overline{AD}^2 - \overline{AF}^2}$.\\
(a) Höhe & 1\E & $h = \sqrt{6.5^2 - 2.5^2} = \sqrt{36} = 6\,\text{cm}$\\
(a) Fläche & 1\F & $A = a \cdot h = 9.5 \cdot 6 = 57\,\text{cm}^2$\\
(b) $h_b$ & 1\F & Dieselbe Fläche mit $b = \overline{BC} = 6.5\,\text{cm}$ als Grundseite: $h_b = \tfrac{A}{b} = \tfrac{57}{6.5} \approx 8.77\,\text{cm}$\\
\end{raster}
\fehler{$h = \sqrt{6.5^2 + 2.5^2} \approx 6.96$ (plus statt minus): Teildreieck- und Höhenzeile 0; $6.96$ ist länger als die Hypotenuse $6.5$ — unmöglich, keine Folgepunkte $\to$ 0 von 4\,P.
Die schräge Seite als Höhe ($h = 6.5$, $A = 61.75$, $h_b = 9.5$): Teildreieck- und Höhenzeile 0; Fläche und $h_b$ als Folgepunkte $\to$ 2 von 4\,P.
$h_b = \tfrac{2A}{b} \approx 17.54$ (Dreiecksformel): $h_b$-Zeile 0.
$h_b = \tfrac{A}{a} = 6$ (durch die falsche Seite geteilt): $h_b$-Zeile 0.}
\end{lpaufgabe}

\begin{lpaufgabe}{G3}{Fliese in Rhombusform}{4 P}
\begin{raster}
(a) Teildreieck & 1 & Die Diagonalen halbieren sich senkrecht; Teildreieck mit der Hypotenuse $a = 8.5$ und der Kathete $\tfrac{e}{2} = 7.5$: $\tfrac{f}{2} = \sqrt{8.5^2 - 7.5^2} = \sqrt{16} = 4\,\text{cm}$\\
(a) Diagonale & 1\E & $f = 2 \cdot 4 = 8\,\text{cm}$\\
(b) Fläche & 1\F & $A = \tfrac{1}{2} \cdot e \cdot f = \tfrac{1}{2} \cdot 15 \cdot 8 = 60\,\text{cm}^2$\\
(c) Abstand & 1\F & Der Rhombus ist ein Parallelogramm: $A = a \cdot h$, also $h = \tfrac{60}{8.5} \approx 7.06\,\text{cm}$\\
\end{raster}
\fehler{$f = 4$ (die halbe Diagonale nicht verdoppelt): Diagonalenzeile 0; $A = 30$ und $h \approx 3.53$ als Folgepunkte $\to$ 3 von 4\,P.
$\tfrac{f}{2} = \sqrt{8.5^2 + 7.5^2} \approx 11.34$ (plus statt minus): Teildreieck- und Diagonalenzeile 0; $11.34$ ist länger als die Hypotenuse $8.5$ — unmöglich, keine Folgepunkte $\to$ 0 von 4\,P.
$A = e \cdot f = 120$ ($\tfrac{1}{2}$ vergessen): Flächenzeile 0; $h = \tfrac{120}{8.5} \approx 14.12$ ist länger als die Seite $8.5$ — unmöglich, Abstandszeile 0 $\to$ 2 von 4\,P.
$h = \tfrac{2A}{a} \approx 14.12$ (Dreiecksformel): Abstandszeile 0 $\to$ 3 von 4\,P.}
\end{lpaufgabe}

\begin{lpaufgabe}{G4}{Walmdach}{4 P}
\begin{raster}
(a) Mittellinie & 1 & $A = m \cdot h$, also $m = \tfrac{A}{h} = \tfrac{38.25}{4.5} = 8.5\,\text{m}$ (gleichwertig: $a + c = \tfrac{2A}{h} = 17\,\text{m}$)\\
(a) First & 1\E & $m = \tfrac{1}{2}(a + c)$, also $c = 2m - a = 17 - 12 = 5\,\text{m}$\\
(b) Überstand & 1 & Teildreieck aus Höhe, Überstand und Grat; Überstand $\tfrac{a - c}{2} = \tfrac{12 - 5}{2} = 3.5\,\text{m}$\\
(b) Grat & 1\F & Der Grat ist die Hypotenuse: $s = \sqrt{3.5^2 + 4.5^2} = \sqrt{32.5} \approx 5.70\,\text{m}$\\
\end{raster}
\fehler{$c = \tfrac{2A}{h} = 17$ ($-\,a$ vergessen): Mittellinienzeile zählt, Firstzeile 0; der Überstand $\tfrac{12 - 17}{2}$ wäre negativ — unmöglich, keine Folgepunkte $\to$ 1 von 4\,P.
$c = \tfrac{A}{h} - a = -3.5$ (Mittellinie als $a + c$ genommen): Mittellinienzeile zählt, die anderen 0 (negative Länge) $\to$ 1 von 4\,P.
Ganzer Unterschied $12 - 5 = 7$ als Überstand ($s = \sqrt{7^2 + 4.5^2} \approx 8.32$): Überstandzeile 0; Gratzeile als Folgepunkt $\to$ 3 von 4\,P.
$s = 3.5 + 4.5 = 8$ (Längen addiert) oder $s = \sqrt{4.5^2 - 3.5^2} \approx 2.83$ (minus statt plus): Gratzeile 0 $\to$ 3 von 4\,P.}
\end{lpaufgabe}

\begin{lpaufgabe}{G5}{Die Mittellinie teilt das Trapez}{4 P}
\begin{raster}
(a) Mittellinie & 1\E & $m = \tfrac{1}{2}(a + c) = \tfrac{1}{2}(11 + 5) = 8\,\text{cm}$\\
(b) Teilhöhen & 1 & Die Mittellinie liegt auf halber Höhe: Beide Teiltrapeze sind $2\,\text{cm}$ hoch; ihre Parallelseiten sind $a$ und $m$ bzw. $m$ und $c$.\\
(b) unten & 1\F & $A_1 = \tfrac{1}{2}(11 + 8) \cdot 2 = 19\,\text{cm}^2$\\
(b) oben & 1\F & $A_2 = \tfrac{1}{2}(8 + 5) \cdot 2 = 13\,\text{cm}^2$ (Probe: $19 + 13 = 32 = m \cdot h$)\\
\end{raster}
\fehler{Beide Teile mit der ganzen Höhe $4$ gerechnet ($A_1 = 38$, $A_2 = 26$): Teilhöhenzeile 0; beide Flächen als Folgepunkte $\to$ 3 von 4\,P.
$m = a + c = 16$ oder $m = \tfrac{a - c}{2} = 3$: Mittellinienzeile 0; die Mittellinie liegt nicht zwischen $5$ und $11$ — unmöglich, keine Folgepunkte für die Flächen; die Teilhöhenzeile zählt $\to$ 1 von 4\,P.
Nur die Gesamtfläche $8 \cdot 4 = 32$: beide Flächenzeilen 0.}
\end{lpaufgabe}

\begin{lpaufgabe}{G6}{Eine fremde Lösung}{4 P}
\begin{raster}
(a) Fehler & 1 & Die Diagonale ist die Hypotenuse, die Höhe eine Kathete — es muss $h^2 = d^2 - b^2$ heissen, nicht plus. Ohne Rechnen: Eine Seite des Rechtecks kann nicht länger sein als seine Diagonale ($85.80 > 65$). Nötig: eine der beiden Begründungen.\\
(b) Höhe & 1\E & $h = \sqrt{65^2 - 56^2} = \sqrt{1089} = 33\,\text{cm}$\\
(b) Umfang & 1\F & $U = 2(56 + 33) = 178\,\text{cm}$\\
(c) Quadrat & 1\E & Im Quadrat ist $d = a\sqrt{2}$, also $a = \tfrac{65}{\sqrt{2}} \approx 45.96\,\text{cm}$ (gleichwertig: $2a^2 = 65^2$)\\
\end{raster}
\fehler{Umfang mit Tims Höhe ($2(56 + 85.80) = 283.60$): Tims Höhe ist länger als die Diagonale — unmöglich, Umfangzeile 0.
$U = 56 + 33 = 89$ (nur zwei Seiten): Umfangzeile 0.
$a = \tfrac{65}{2} = 32.5$ (Diagonale halbiert) oder $a = 65\sqrt{2} \approx 91.92$: Quadratzeile 0.}
\end{lpaufgabe}

\needspace{14\baselineskip}
\section*{4 · Selbsteinschätzung}
\begin{tabularx}{\linewidth}{@{}l X@{}}\toprule
21 – 24 P & Die geprüften Teile sitzen. Wo du Punkte verloren hast: das Kapitel dieser Aufgabe nochmals (Zuordnung unten).\\
16 – 20 P & Das schwächste Kapitel nochmals: Tüfteln und Übungen des Kapitels, in dem du die meisten Punkte verloren hast.\\
11 – 15 P & Zurück zu den Kapiteln aller Aufgaben, in denen du Punkte verloren hast.\\
0 – 10 P & Zurück zu Kapitel 1 und von dort der Reihe nach weiter.\\\midrule
\multicolumn{2}{@{}p{\linewidth}@{}}{\small\textbf{Aufgabe $\to$ Kapitel:} G1 $\to$ Kapitel 1 (Die Vierecks-Familie); G2 $\to$ Kapitel 2 (Fläche und Umfang), die Höhe in Kapitel 4; G3 $\to$ Kapitel 4 (Fehlende Längen), Fläche und Abstand in Kapitel 2; G4 $\to$ Kapitel 3 (Trapez und Mittellinie), der Grat in Kapitel 4; G5 $\to$ Kapitel 3; G6 $\to$ Kapitel 4, der Umfang in Kapitel 2}\\\bottomrule
\end{tabularx}
\end{document}
